\documentclass[10pt,journal,compsoc]{IEEEtran}
\usepackage[utf8]{inputenc}
\usepackage[T1]{fontenc}
\usepackage{url}
\usepackage[hidelinks]{hyperref}
\usepackage{cite}
\usepackage{amsmath,amssymb,amsfonts}
\usepackage{algorithmic}
\usepackage{graphicx}
\usepackage{textcomp}
\usepackage{xcolor}
\usepackage{tabularx}
\usepackage{listings}
\usepackage{array}

\lstset{
  basicstyle=\small\ttfamily,
  breaklines=true,
  frame=single,
  columns=fullflexible
}

\begin{document}

\title{Novel Contributions}
\author{Bryan Alexander Buchorn \\ Independent Researcher \\ bryan.alexander@buchorn.com}
\markboth{CEREBRUM v1.2.4 Hardened Edition · March 2026}{}

\maketitle

\begin{abstract}
This module forms a critical component of the CEREBRUM framework, a community-structured graph attention engine for Knowledge Graph reasoning. It focuses on Framework Implementation Guide.
\end{abstract}

\begin{IEEEkeywords}
Knowledge Graph, Graph Attention, DSCF, CSA, Hierarchical Reasoning.
\end{IEEEkeywords}

\section{Novel Contri\-butions and IP Claims}

\subsection{CEREBRUM: Community-Structured Graph Attention for Knowledge Graph Reasoning}

\textbf{Document Classi\-fication}: Intell\-ectual Property Reference
\textbf{Authors}: Bryan Alexander Buchorn
\textbf{Date}: March 2026
\textbf{Version}: v1.1.0 (Phase 20 COMPLETE)

\textgreater  This document conso\-lidates the novel technical contr\-ibutions of the CEREBRUM framework for use in patent appli\-cations, academic priority claims, and commercial IP protection. Each claim is subst\-antiated with prior art analysis and a statement of the specific technical distinction.

---

\subsection{Part I: Core Algorithm Claims}

\subsubsection{Claim 1: Dual/Triple Signal Community Fusion (DSCF / TSC)}

\textbf{Description}: A community detection algorithm that applies Local Propagation Algorithm (LPA), modularity gain (Louvain-style), and optionally Infomap flow signals \textbf{simul\-taneously at each node update}, fusing them via a weighted consensus mechanism at each iteration.

\textbf{Novelty Statement}: All prior community detection algorithms that combine multiple signals do so by operating on disjoint node subsets (e.g., apply LPA to low-degree nodes, Louvain to high-degree nodes [Sun et al., 2024]) or by running algorithms seque\-ntially and merging results. DSCF is the first algorithm that applies all signals to every node at every update step. The simul\-taneous per-node fusion produces communities with distinct "dual-signal" structural character that is not achievable by sequential or partitioned combination.

\textbf{Closest Prior Art}:
\begin{itemize}
\item Louvain (Blondel et al., 2008): Global modularity only, no LPA signal
\item Leiden (Traag et al., 2019): Refinement phase + modularity, no LPA signal
\item LPA (Raghavan et al., 2007): Local propagation only, no modularity signal
\item LPA-Louvain hybrids (Sun et al., 2024): Apply signals to disjoint node subsets by degree threshold — categ\-orically different from DSCF's per-node simul\-taneous fusion

\end{itemize}
\textbf{Key Technical Differ\-entiator}: The specific archi\-tectural choice to compute both $\Delta Q_{modularity}(v, c)$ and $f_{LPA}(v, c)$ for every candidate community $c$ at every node $v$ in every iteration, then fuse them via: $\text{score}(v, c) = \alpha \cdot \Delta Q(v,c) + \beta \cdot f_{LPA}(v,c) + \gamma \cdot \text{flow}(v,c)$

\textbf{Relevant files}: \texttt{core/community\_engine.py}, \texttt{core/leiden\_native.py}
\textbf{Documented in}: \texttt{docs/arxiv/PAPER\_001\_DSCF\_TSC.md}, \texttt{docs/specifications/SPEC\_001\_DSCF\_TSC.md}

---

\subsubsection{Claim 2: Community-Structured Attention (CSA) Formula}

\textbf{Description}: A graph edge attention weight formula that incor\-porates community membership as a soft global constraint alongside semantic similarity, relation type, path length penalty, hop decay, and PageRank centrality. The formula is training-free, computed analy\-tically from graph topology at query time.

\textbf{The Formula}:
\begin{equation}
a(u,v,k) = \sigma\left(\alpha \cdot \cos(\vec{e}_u, \vec{e}_v) + \beta \cdot S_C(u,v) + \gamma \cdot w_{rel} - \delta \cdot d_{norm}(u,v) + \varepsilon \cdot \phi(k) + \zeta \cdot PR(v)\right)
\end{equation}

\textbf{Novelty Statement}: Graph Attention Networks (GAT, HAN, HGT) compute local attention over immediate neigh\-borhoods using learned weight matrices. CSA is the first attention formulation that includes global community membership ($S_C(u,v)$) as a term. No published GNN or KG reasoning paper computes attention weights using community assignment as an explicit term in the attention formula.

\textbf{Closest Prior Art}:
\begin{itemize}
\item GAT (Veličković et al., 2018): Local neigh\-borhood attention, learned weights, no community term
\item HAN (Wang et al., 2019): Meta-path attention, no community term
\item HGT (Hu et al., 2020): Hetero\-geneous attention, learned, no community term
\item GraphRAG (Microsoft, Edge et al., 2024): Communities used for LLM summa\-rization, not as attention weights in traversal

\end{itemize}
\textbf{Key Technical Differ\-entiator}: The $\beta \cdot S_C(u,v)$ term, where $S_C(u,v) \in [0,1]$ measures community co-membership and is computed from live DSCF partitions, not from learned parameters.

\textbf{Relevant files}: \texttt{core/attention\_engine.py}
\textbf{Documented in}: \texttt{docs/arxiv/PAPER\_002\_CSA.md}, \texttt{docs/specifications/SPEC\_002\_CSA.md}

---

\subsubsection{Claim 3: CSA-Guided Multi-Hop Beam Search (Glass-Box Reasoning)}

\textbf{Description}: A multi-hop KG traversal archi\-tecture where beam search is guided exclusively by CSA attention weights derived from graph topology, producing a complete verifiable reasoning trace at zero inference cost (no trained model, no LLM).

\textbf{Novelty Statement}: All published multi-hop KG reasoning systems using beam search (MINERVA, DeepPath, BeamQA) require training on labeled question-answer pairs. CEREBRUM's BeamTraversal uses CSA weights computed analy\-tically from graph structure — achieving competitive H@10 on MetaQA zero-shot, with no training data.

\textbf{Closest Prior Art}:
\begin{itemize}
\item MINERVA (Das et al., 2018): RL-trained path agent, requires labeled QA pairs
\item DeepPath (Xiong et al., 2017): RL, requires training
\item BeamQA (Atif et al., 2023): Supervised sequence-to-sequence beam search

\end{itemize}
\textbf{Key Technical Differ\-entiator}: Zero-shot multi-hop reasoning at competitive H@10 using only graph topology.

\textbf{Relevant files}: \texttt{reasoning/traversal.py}, \texttt{reasoning/path\_scorer.py}

---

\subsection{Part II: Neurom\-orphic Archit\-ecture Claims}

\subsubsection{Claim 4: Experience-Dependent Bridge Twin Nodes (LTP/LTD Analog for KG)}

\textbf{Description}: A mechanism that mater\-ializes structural relay nodes (bridge twins) between graph communities based on traversal co-occurrence frequency, analogous to synaptic LTP/LTD in hippocampal circuits. Bridge twins form when crossing count exceeds a threshold and are pruned when crossings decay below a threshold.

\textbf{Novelty Statement}: No published KG system uses experience-dependent structural modif\-ication inspired by LTP/LTD. Static graph shortcuts (agentic LLM expansion, shortcut edges in GNNs) are added by external agents or hand-coded rules, not by the reasoning engine's own traversal history. The specific mechanism of counting traversal co-occurrences and triggering structural node mater\-ialization from those counts is, to our knowledge, original.

\textbf{Relevant files}: \texttt{core/bridge\_engine.py}
\textbf{Documented in}: \texttt{docs/arxiv/PAPER\_003\_BRIDGE\_TWINS.md}

---

\subsubsection{Claim 5: STDP-Inspired Directional Causal Edge Inference}

\textbf{Description}: A streaming discretizer that mater\-ializes directional \texttt{CAUSES} edges in a Knowledge Graph from event timing patterns, using Hebbian-inspired weight accum\-ulation with confi\-gurable LTP/LTD asymmetry and lazy O(1) weight decay.

\textbf{Novelty Statement}: STDP has been applied to spiking neural network learning rules extensively, but no published work applies STDP mechanics to dynamic causal edge discovery in Knowledge Graphs. The combination of: (1) Hebbian co-occurrence weighting, (2) directional edge mater\-ialization (pre$\to$post), (3) lazy weight decay reducing per-event complexity from O(N) to O(1), and (4) adversarial causal flood protection (min\_causal\_span, chi-squared) is an original contr\-ibution.

\textbf{Relevant files}: \texttt{core/discretizer.py}
\textbf{Documented in}: \texttt{docs/arxiv/PAPER\_004\_STDP\_CAUSAL.md}

---

\subsubsection{Claim 6: REM Cycle — Sleep-Inspired KG Maintenance}

\textbf{Description}: A three-phase background maintenance loop (Prune $\to$ Consolidate $\to$ Synthesize) for autonomous KG quality management, inspired by the biological sleep-cycle memory conso\-lidation model (synaptic homeostasis, NREM slow-wave, hippocampal replay).

\textbf{Novelty Statement}: KG maintenance systems (NELL, ATOMIC) focus on \textit{adding} new facts. The REM Cycle's three-phase design for autonomous removal, structural conso\-lidation, and proactive synthesis with diffe\-rential decay rates per edge relation type has no precedent in the KG maintenance literature.

\textbf{Relevant files}: \texttt{core/rem\_engine.py}
\textbf{Documented in}: \texttt{docs/arxiv/PAPER\_007\_REM\_CYCLE.md}

---

\subsection{Part III: Cross-Modal and Federated Claims}

\subsubsection{Claim 7: Procrustes SVD Sensor-to-KG Alignment with Canonical Basis Anchor}

\textbf{Description}: A cross-modal alignment system that projects heter\-ogeneous sensor signals (time-series, waveforms) into the same embedding space as Knowledge Graph text entities via Procrustes SVD rotation, with a canonical embedding basis that prevents geometric drift accum\-ulation across federated hops.

\textbf{Novelty Statement}: Multimodal KG papers (MMKG, MKGFormer) integrate images and text but not continuous sensor signals. The specific combination of: (1) statistical/spectral signal feature extraction, (2) Procrustes SVD alignment to KG embedding space, (3) a canonical root basis shared across all federation parti\-cipants, and (4) namespace isolation preventing cross-modal entity collisions, is original.

\textbf{Relevant files}: \texttt{core/signal\_encoder.py}
\textbf{Documented in}: \texttt{docs/arxiv/PAPER\_008\_SIGNAL\_ENCODER.md}

---

\subsubsection{Claim 8: Holographic Index for Privacy-Preserving Federated Discovery}

\textbf{Description}: A federated graph discovery system using Bloom filter sketches (proba\-bilistic membership) and community centroid vectors (structural finge\-rprints) to enable two CEREBRUM nodes to identify likely-relevant shared content without revealing individual entity identities.

\textbf{Novelty Statement}: Federated learning (McMahan et al., 2017) aggregates model gradients, not graph structure. SPARQL federation requires full query exposure. The Holographic Index is the first KG-specific privacy-preserving discovery mechanism using the combination of Bloom filters and community centroid structural finge\-rprints.

\textbf{Relevant files}: \texttt{core/hologram.py}
\textbf{Documented in}: \texttt{docs/arxiv/PAPER\_005\_HOLOGRAPHIC\_INDEXING.md}

---

\subsection{Part IV: Production Hardening Claims}

\subsubsection{Claim 9: Bayesian Beam Search with Warm-Start Prior Seeding for KG Traversal}

\textbf{Description}: Application of Beta-distr\-ibution path modeling and Thompson sampling to multi-hop KG beam search, with a warm-start mechanism that seeds the first-hop Beta prior from the CSA attention weight to reduce cold-start variance.

\textbf{Novelty Statement}: Bayesian bandits and Thompson sampling are well-established in recom\-mendation and exploration-explo\-itation literature. Their application to \textit{multi-hop KG beam search with path-level Beta distr\-ibutions} seeded from graph-structural attention weights is novel — no prior KG reasoning paper uses this formulation.

\textbf{Relevant files}: \texttt{reasoning/traversal.py}
\textbf{Documented in}: \texttt{docs/arxiv/PAPER\_006\_BAYESIAN\_BEAM.md}

---

\subsubsection{Claim 10: Eight-Hole Structural Hardening Taxonomy for KG Systems}

\textbf{Description}: A systematic methodology for identifying and patching cross-feature interaction bugs ("structural holes") in complex KG systems, covering five taxonomic categories: stale reference, adversarial input, namespace collision, bias/saturation, and evaluation artifacts.

\textbf{Novelty Statement}: The specific taxonomy of structural hole categories as applied to KG reasoning systems — and the complete set of eight patched holes with backward-compatible fixes — constitutes a novel engineering contr\-ibution for production KG deployment.

\textbf{Documented in}: \texttt{docs/arxiv/PAPER\_016\_PRODUCTION\_HARDENING.md}

---

\subsection{Part V: Composite System Claims}

\subsubsection{Claim 11: CEREBRUM as an Integrated Zero-Shot KG Reasoning System}

\textbf{Description}: The complete CEREBRUM system integrating DSCF community detection, CSA attention, zero-shot beam traversal, LTP/LTD bridge twins, STDP causal inference, REM Cycle maintenance, Bayesian beam search, cross-modal alignment, federated discovery, streaming ingest, metac\-ognitive verif\-ication, and production hardening into a single coherent framework with zero training data requi\-rements and complete reasoning trans\-parency.

\textbf{Novelty Statement}: No prior work integrates this combination of capab\-ilities in a single KG reasoning framework. The specific archi\-tectural choices — using communities as attention heads, using traversal co-occurrence for structural modif\-ication, using STDP timing for causal discovery, and composing these into a sleep-cycle maintenance loop — represent a novel systems-level contr\-ibution.

---

\subsection{Prior Art Summary Table}

\begin{center}
{\footnotesize
\begin{tabularx}{\columnwidth}{|>{\raggedright\arraybackslash}X|>{\raggedright\arraybackslash}X|>{\raggedright\arraybackslash}X|}
\hline
CEREBRUM Component & Closest Prior Art & Key Distinction \\ \hline
DSCF simul\-taneous fusion & LPA-Louvain hybrids (Sun 2024) & Per-node vs. per-population fusion \\ \hline
CSA formula ($S_C$ term) & GAT, HAN, HGT & Community term absent from all GNN attention \\ \hline
Zero-shot beam traversal & MINERVA, DeepPath, BeamQA & Training required vs. fully training-free \\ \hline
Bridge Twins (LTP/LTD) & GNN shortcuts, agentic expansion & Experience-dependent vs. static/agent-added \\ \hline
STDP causal edges in KG & SNN-STDP & Neural learning rule vs. KG edge discovery \\ \hline
REM Cycle maintenance & NELL, ATOMIC, KGE training & Proactive pruning + synthesis vs. fact addition \\ \hline
Procrustes sensor$\to$KG & MMKG, MKGFormer & Sensor signals vs. image/text only \\ \hline
Holographic Index & Federated learning, SPARQL federation & Privacy-preserving structural discovery \\ \hline
Bayesian beam + warm-start & Bandit algorithms & KG-specific Beta seeding from CSA weights \\ \hline
Structural hole taxonomy & Standard testing methodology & Cross-feature interaction analysis for KG \\ \hline
\end{tabularx}
}
\end{center}

---

\subsection{Legal Notice}

All rights, title, and interest in and to the CEREBRUM software, docum\-entation, algorithms, and related intel\-lectual property documented herein are and shall remain the exclusive property of \textbf{Bryan Alexander Buchorn}.

CEREBRUM is dual-licensed:
\begin{itemize}
\item \textbf{Non-Commercial Use}: PolyForm Noncom\-mercial License 1.0.0
\item \textbf{Commercial Use}: Separate commercial license required

\end{itemize}
For commercial licensing: \textbf{bryan.alexander@buchorn.com}

\textbf{Copyright © 2026 Bryan Alexander Buchorn. All Rights Reserved.}


\bibliographystyle{IEEEtran}
\bibliography{../../docs/latex/references}

\end{document}
